\chapter{Introdução}
\label{cap:introducao}

Água e energia representam duas das maiores despesas da sociedade moderna. Globalmente, estima-se que 2-3\% do consumo energético é utilizado para o bombeamento e tratamento de água em contextos urbanos e industriais \cite{M2012}. Aproximadamente 86,4\% dos custos gerados por uma bomba de água estão relacionados à energia consumida durante sua operação no abastecimento \cite{Nault2015}, e estudos apontam que a eficiência energética desses sistemas pode ser melhorada em pelo menos 25\% \cite{Moreira2013}.

Diante da crescente demanda e das preocupações ambientais, otimizar o uso desses recursos é imperativo para garantir operações sustentáveis e eficientes. Segundo um relatório da Agência Internacional de Energia \cite{ieaWEO2016Special}, o consumo energético no setor de água mais que duplicará nos próximos 20 anos, principalmente por conta de projetos de dessalinização.

A gestão do consumo energético em sistemas de distribuição de água (SDAs) é uma questão multifacetada. Infraestruturas antigas, rotinas operacionais ineficientes e equipamentos obsoletos são apenas alguns dos obstáculos \cite{M2012, moura2010relaccao}. No entanto, novas tecnologias, como conversores de frequência, algoritmos de otimização e sistemas de monitoramento em tempo real com IoT, surgem como soluções promissoras \cite{Mackle1995, moura2010relaccao, ufopREPOSITORIOINSTITUCIONAL}.

Este trabalho aborda a otimização operacional de bombas em SDAs, um problema desafiador devido à natureza discreta do espaço de busca, não convexidade das leis físicas envolvidas, necessidade de simulações hidráulicas, flutuações nas demandas e complexidade para quantificar as despesas operacionais \cite{Makaremi2017}.

Existem duas principais abordagens para esse problema: a modernização de equipamentos, que demandam investimentos significativos, e a otimização da programação operacional das bombas, que pode ser implementada com baixo custo, através da capacitação dos operadores e adoção de sistemas de apoio a tomada de decisão \cite{Moreira2013}. Esta última pode oferecer economias imediatas nos custos energéticos, destacando-se na literatura por sua viabilidade.

Os custos de operação de SDAs, por sua vez, compreendem tanto despesas de consumo elétrico, como despesas relacionadas à manutenção das bombas de água \cite{vanZyl2004}. A manutenção, apesar de ser um fator de custo significativo, apresenta dificuldades de quantificação devido a várias influências, como o tipo e idade da bomba, o perfil de operação, a experiência do operador e as condições climáticas \cite{Makaremi2017}.

Neste artigo, seguindo as diretrizes de amplos trabalhos da literatura \cite{lopez2008ant, lansey1994optimal, boulos2001optimal, vanZyl2004, savic1997multiobjective, lopez2005multi, kelner2003optimal, bene2013comparison}, assumimos o desgaste através de uma medida substituta, o número de acionamentos de bombas.

Otimizar a operação de SDAs é um campo de pesquisa ativo por quase meio século \cite{MalaJetmarova2017}. Inicialmente, se utilizavam modelos hidráulicos simplificados, funções de custo baseadas em tarifas fixas \cite{MalaJetmarova2017} e a abordagem se dava através de métodos determinísticos como Programação Linear \cite{jowitt1992optimal}, Programação Não Linear \cite{chase1993computer} e Programação Dinâmica \cite{Nitivattananon1996, sterling1975dynamic, zessler1989optimal}.

Desde a década de 1990, algoritmos metaheurísticos, como Algoritmos Genéticos (AGs) \cite{Mackle1995, vanZyl2004, savic1997multiobjective}, Recozimento Simulado (SA) \cite{goldman1999application,teegavarapu2002optimal}, Otimização por Colônia de Formigas (ACO) \cite{lopez2008ant} e Enxame de Partículas \cite{ostadrahimi2012multi}, emergiram como alternativas promissoras para lidar com a complexidade do problema de programação de bombas \cite{MalaJetmarova2017}. Embora populares, esses métodos não garantem a solução ótima global e podem enfrentar dificuldades na busca por soluções de alta qualidade, principalmente em problemas com muitas restrições e redes mais complexas \cite{Costa2015}.

A crescente complexidade dos SDAs e a necessidade de considerar um maior número de variáveis e restrições motivaram a busca por métodos mais eficientes. A estrutura tarifária da eletricidade, por exemplo, é um fator crucial na otimização dos custos de bombeamento, exigindo estratégias que maximizem o uso em horários de menor tarifa.

Avanços tecnológicos permitiram a incorporação de simuladores hidráulicos mais realistas, como o EPANET \cite{epanet}, nos métodos de otimização. Contudo, o esforço computacional necessário para tais simulações, principalmente em modelos com redes mais complexas, ainda limitam sua aplicabilidade \cite{MalaJetmarova2017}.

A combinação de métodos determinísticos e metaheurísticas se mostraram uma estratégia eficaz para melhorar o desempenho dos algoritmos e reduzir o tempo de processamento. Por exemplo, a hibridização de AGs com métodos de busca local puderam acelerar a convergência e melhorar a qualidade das soluções \cite{reis2006water, vanZyl2004}.

No entanto, estudos mostram que essas metaheurísticas, quando integradas a simuladores de rede como o EPANET, podem limitar sua implementação em grandes SDAs e em cenários de controle em tempo real, devido ao elevado esforço computacional necessário \cite{MalaJetmarova2017}. Em resposta, métodos determinísticos mais eficientes voltaram a ter destaque e têm sido cada vez mais aplicados.

Em muitos cenários, o número de bombas em SDAs não é excessivamente alto, tornando viável a aplicação de métodos exatos que exploraram a enumeração das soluções possíveis, uma vez que o número de combinações é computacionalmente tratável \cite{Costa2015}. Algoritmos determinísticos, como o Branch-and-Bound (B\&B) \cite{land2010automatic, costa2001global}, além de métodos como Decomposições Lagrangianas \cite{Ghaddar2015} e de Benders generalizadas \cite{verleye2016generalized}, têm se mostrado eficientes na otimização da programação de bombas, ao dividirem problemas complexos em subproblemas menores e mais gerenciáveis.

Apesar dos progressos, desafios permanecem na otimização da programação de bombas, incluindo escalabilidade, transferibilidade do modelo (capacidade do modelo de ser aplicado em diferentes redes), incorporação de incertezas na demanda de água e o alto custo computacional \cite{Kerimov2023}. O uso de metamodelos, como Redes Neurais Artificiais (RNAs)\cite{behzadian2009stochastic, broad2005water, rao2007use}, para simular o comportamento dos SDAs, podem superar alguns desses desafios, reduzindo tempos de simulação e permitindo a implementação de estratégias de controle em tempo real.

Não há consenso sobre qual método é o mais adequado para a otimização da operação de SDAs, pois nenhum se mostra totalmente satisfatório ao enfrentar os diversos desafios envolvidos \cite{MalaJetmarova2017}. Nesse contexto, o presente artigo propõe um algoritmo Branch-and-Bound (B\&B) como método exato para encontrar a solução ótima global para a programação de bombas. 

O B\&B utiliza uma abordagem de divisão e conquista, dividindo o problema em subproblemas menores e explorando o espaço de soluções de maneira sistemática. Ao fixar variáveis de decisão e empregar limites inferiores, o algoritmo elimina soluções inviáveis, reduzindo significativamente o espaço de busca e permitindo encontrar soluções ótimas globais \cite{Costa2015}.

Este trabalho é uma extensão do estudo de Costa \cite{Costa2015}, que apresentou algumas limitações, especialmente em relação ao elevado tempo computacional, causado pelas inúmeras simulações hidráulicas no EPANET e pela ausência de uma estratégia eficaz para o cálculo de limites inferiores, o que dificultou a poda eficiente de soluções inviáveis. Além disso, a busca em largura (Breadth-First Search) exigia maior uso de memória e processamento, pois todas as soluções parciais de cada nível eram analisadas antes de avançar na árvore.

A nova abordagem proposta neste trabalho supera essas limitações ao introduzir uma busca em profundidade (Depth-First Search) especializada. Diferente da busca em largura, que não calculava limites inferiores, esta implementação atualiza continuamente o limite inferior conforme o algoritmo avança. Isso é feito através do cálculo de um peso associado a cada decisão, normalizado pelo nível dos nós da árvore, o que permite priorizar nós mais profundos. Com essa estratégia, o algoritmo faz o backtracking de forma mais eficiente e percorre a árvore de maneira aprimorada, reduzindo o tempo computacional na busca por soluções ótimas globais. Além disso, a aplicação de técnicas de paralelismo melhora ainda mais o desempenho, tornando o algoritmo mais eficiente e adequado para redes de maior escala e complexidade.
